\documentclass[12pt]{extarticle}

% Unified preamble from MASTER_COMBINED.tex
\usepackage{amsmath}
\usepackage{amssymb}
\usepackage{amsthm}
\usepackage{graphicx}
\usepackage{xcolor}
\usepackage[most]{tcolorbox}
\usepackage{tikz}
\usepackage{fancyhdr}
\usepackage{geometry}
\usepackage{array}
\usepackage{booktabs}
\usepackage{listings}
\usepackage{enumitem}
\usepackage{cancel}
\usepackage{setspace}
\usepackage[hidelinks]{hyperref}
\usepackage{tikzpagenodes}

\geometry{margin=1in}

\setlength{\parskip}{0.45em}
\setlength{\parindent}{0pt}
\setlength{\headheight}{15pt}
\setstretch{1.18}
\raggedbottom

% Global header: © 2025 Pranav Ramesh. All rights reserved. www.lambdamath.dev on every page
\pagestyle{fancy}
\fancyhf{}
\fancyhead[C]{\textcopyright\ 2025 Pranav Ramesh. All rights reserved. www.lambdamath.dev}
\fancyfoot[C]{\thepage}\fancyfoot[R]{\small\textcopyright\ 2025 Pranav Ramesh.}
\renewcommand{\headrulewidth}{0.4pt}
\renewcommand{\footrulewidth}{0pt}

\fancypagestyle{plain}{
    \fancyhf{}
    \fancyhead[C]{\textcopyright\ 2025 Pranav Ramesh. All rights reserved. www.lambdamath.dev}
    \fancyfoot[C]{\thepage}\fancyfoot[R]{\small\textcopyright\ 2025 Pranav Ramesh.}
    \renewcommand{\headrulewidth}{0.4pt}
    \renewcommand{\footrulewidth}{0pt}
}

\usetikzlibrary{shapes, arrows, positioning, calc, angles, quotes}

% Watermark
\usepackage{background}
\usepackage{hyperref} % if already loaded elsewhere, keep it
\usepackage{tikz}

\backgroundsetup{
    scale=1, % Increased scale to cover more area
    angle=0,
    opacity=0.1,
    contents={
        \tikz[remember picture, overlay]
            \node[
                rotate=45,
                font=\fontsize{300}{300}\bfseries, % Increased font size
                inner sep=0pt,
                color=gray % Changed color to gray
            ]
            at (current page.center)
            {LambdaMath.dev};
    }
}

% Unified styled boxes
\newtcolorbox{conceptbox}[1][]{
    colback=blue!5!white,
    colframe=blue!75!black,
    title=#1,
    fonttitle=\bfseries
}

\newtcolorbox{remarkbox}{
    colback=green!5!white,
    colframe=green!60!black,
    fonttitle=\bfseries,
    title=Remark
}

\newtcolorbox{examplebox}{
    colback=orange!5!white,
    colframe=orange!75!black,
    fonttitle=\bfseries,
    title=Example
}

\newtcolorbox{warningbox}{
    colback=red!5!white,
    colframe=red!75!black,
    fonttitle=\bfseries,
    title=Warning
}

\newtcolorbox{theorembox}{
    colback=purple!5!white,
    colframe=purple!75!black,
    fonttitle=\bfseries,
    title=Theorem
}

\newtcolorbox{solutionbox}{
    colback=yellow!5!white,
    colframe=orange!75!black,
    fonttitle=\bfseries,
    title=Solution
}

\newtcolorbox{insightbox}{
    colback=teal!5!white,
    colframe=teal!70!black,
    fonttitle=\bfseries,
    title=Insight / Remark
}

\newtcolorbox{methodsummarybox}{
    colback=gray!6!white,
    colframe=gray!70!black,
    fonttitle=\bfseries,
    title=Method Summary
}

\newtcolorbox{commonpitfallsbox}{
    colback=red!3!white,
    colframe=red!60!black,
    fonttitle=\bfseries,
    title=Common Pitfalls
}

% Difficulty tag and labels
\newcommand{\difftag}[1]{\textbf{[#1]}}
\newcommand{\difficulty}[1]{\textbf{(#1)}\,}

% Proof-style environments
\newenvironment{FormalProof}{\begin{solutionbox}\textbf{Formal Proof.} }{\end{solutionbox}}
\newenvironment{ProofSketch}{\begin{solutionbox}\textbf{Proof Sketch.} }{\end{solutionbox}}
\newenvironment{Heuristic}{\begin{insightbox}\textbf{Heuristic / Intuition.} }{\end{insightbox}}

% Numbered theorems within sections
\newtcbtheorem[number within=section]{definition}{Definition}{
    colback=gray!5!white,
    colframe=gray!50!black,
    coltitle=black,
    fonttitle=\bfseries,
    title={Definition~\thetcbcounter}
}{def}

\newtcbtheorem[number within=section]{theorem}{Theorem}{
    colback=purple!5!white,
    colframe=purple!70!black,
    coltitle=black,
    fonttitle=\bfseries,
    title={Theorem~\thetcbcounter}
}{thm}

\newtcbtheorem[number within=section]{strategy}{Strategy}{
    colback=teal!4!white,
    colframe=teal!70!black,
    coltitle=black,
    fonttitle=\bfseries,
    title={Strategy~\thetcbcounter}
}{strat}

\newtcbtheorem[number within=section]{example}{Example}{
    colback=orange!5!white,
    colframe=orange!75!black,
    coltitle=black,
    fonttitle=\bfseries,
    title={Example~\thetcbcounter}
}{ex}

\newtcbtheorem[number within=section]{commonmistake}{Common Mistake}{
    colback=red!5!white,
    colframe=red!75!black,
    coltitle=black,
    fonttitle=\bfseries,
    title={Common Mistake~\thetcbcounter}
}{cm}

\begin{document}

\begin{titlepage}
\centering
\vspace*{2cm}

{\Huge\bfseries Number Theory for AMC}\\[0.5cm]
{\Large Competition Problem Solving}\\[2cm]

{\large AMC 10 $\cdot$ AMC 12 $\cdot$ AIME $\cdot$ IOQM}\\[3cm]

{\Large A Strategic Guide to\\[0.3cm]
Structure, Invariants, and Modular Thinking}\\[3cm]

{\Large\textbf{Pranav Ramesh}}\\[3cm]

\vfill

{\large 23 Dec 2025 (Version 1)}
\end{titlepage}

% Copyright Page
\thispagestyle{empty}
\vspace*{\fill}
\begin{center}
\textbf{\Large Copyright \textcopyright\ 2025 Pranav Ramesh}

All rights reserved.

This work is protected under the Copyright Act, 1957 (India).

\vspace{1.5cm}

This publication is the original work of Pranav Ramesh published under the imprint LambdaMath.

\vspace{1.5cm}

No part of this work may be reproduced, stored, transmitted, or distributed in any form or by any means, electronic, mechanical, photocopying, recording, or otherwise, without prior written permission of the copyright holder, except for brief quotations for educational or review purposes.

\vspace{1.5cm}

First published on 23 Dec 2025 at: www{.}lambdamath{.}dev

For permissions: pranavramesh2017@gmail{.}com
\end{center}
\vspace*{\fill}

\newpage

% Dedication Page
\thispagestyle{empty}
\vspace*{\fill}
\begin{center}
\Large

\textbf{Dedicated to}

\vspace{1cm}

Mr. Krishna Rao, My highschool mathematics teacher,

Mr. Siva Rama Praveen K, Principal,

Mr. Rajashekar Reddy, AGM, and

Mrs. Sumathi, My mathematics mentor,

\vspace{1cm}

for their guidance, encouragement, and belief in the value of rigorous thinking.
\end{center}
\vspace*{\fill}

\newpage

	\tableofcontents
\newpage

\section*{Preface}
\addcontentsline{toc}{section}{Preface}

\subsection*{Who This Book Is For}
This book is designed for students preparing for AMC 10/12, AIME, and IOQM who want a competition-focused approach to number theory.

\textbf{You should use this book if you:}
\begin{itemize}
    \item Want to recognize number-theoretic structures quickly (primes, orders, residues)
    \item Need a systematic approach to modular arithmetic and factorization
    \item Prefer understanding \emph{why} tools work, not just memorizing statements
\end{itemize}

\subsection*{What Makes This Book Different}
Competition number theory values structure and inevitability over brute force. We emphasize invariants, congruences, and strategic prime factor thinking so you can see the path before calculating.

\subsection*{How to Use This Book}
\begin{enumerate}
    \item \textbf{Intuition first:} Read concept boxes before formulas.
    \item \textbf{Tagging:} (M) must memorize; (R) recognize and know when to apply.
    \item \textbf{Practice:} Attempt examples before reading solutions.
    \item \textbf{Review:} Revisit core tools (orders, LTE, CRT) until they are automatic.
\end{enumerate}

\subsection*{Colored Boxes Guide}
\begin{itemize}
    \item \textcolor{blue!75!black}{\textbf{Concepts:}} Core ideas and methods
    \item \textcolor{orange!75!black}{\textbf{Examples:}} Worked problems with detailed solutions
    \item \textcolor{green!60!black}{\textbf{Remarks:}} Strategic insights and tips
    \item \textcolor{red!75!black}{\textbf{Warnings:}} Common mistakes to avoid
    \item \textcolor{purple!75!black}{\textbf{Theorems:}} Formal statements you should know
\end{itemize}

\subsection*{Study Recommendations}
\begin{itemize}
    \item Read actively with paper and pencil; test small numeric cases
    \item Try examples yourself before reading solutions
    \item Memorize key facts (orders, totient values, small mod cycles)
    \item Return to review sections regularly—mastery requires repetition
\end{itemize}

\subsection*{Prerequisites}
Comfort with algebra, basic modular arithmetic, and simple proof techniques (contradiction, induction).

\subsection*{Beyond This Book}
Pair these notes with past AMC/AIME problems. After solving, ask: "Which invariant or congruence unlocked this?" Keep a running list of patterns and traps.

\subsection*{Acknowledgements}
Thanks to the problem authors and the math contest community whose ideas shape this material.

\newpage

\section*{Chapter 2: Number Theory (AMC 12 / AIME / IOQM Focus)}

\subsection*{Prerequisites}
Comfort with algebra, basic modular arithmetic, and simple proof techniques (contradiction, induction).

\subsection*{Voice and Style}
We aim for an "AMC proof style": concise, logical, and pattern-driven. Remarks highlight strategy; warnings flag common traps.

\newpage

\section*{Chapter 2: Number Theory (AMC 12 / AIME / IOQM Focus)}

This chapter develops \textbf{high-level number theory tools} required for strong AMC 12, AIME, and IOQM performance. The focus is on \textbf{structure, invariants, and inevitability} rather than computation. Every topic is explained with intent: \emph{why it exists}, \emph{when it is used}, and \emph{how it interacts with other tools}.

\begin{conceptbox}[Philosophy of Competition Number Theory]
Number theory problems require:
\begin{enumerate}
    \item \textbf{Prime Decomposition:} Breaking numbers into prime factors
    \item \textbf{Modular Thinking:} Working with remainders and cycles
    \item \textbf{Divisibility Logic:} Understanding when and why division works
    \item \textbf{Strategic Calculation:} Computing only what's necessary
\end{enumerate}
\end{conceptbox}

Throughout this chapter:
\begin{itemize}
\item \textbf{(M)} = Must-memorize (instant recall)
\item \textbf{(R)} = Recognition-based (know when and why it applies)
\end{itemize}

\section{Primes and Prime Factorization}

Primes are the \textbf{atoms of arithmetic}. Almost every AMC-level number theory problem ultimately reduces to prime behavior.

\subsection*{Fundamental Theorem of Arithmetic}

\begin{theorembox}
\textbf{(M)} Every integer $n>1$ can be written uniquely as
\[
n = p_1^{a_1} p_2^{a_2} \cdots p_k^{a_k},
\]
where the $p_i$ are distinct primes and $a_i > 0$.
\end{theorembox}

\begin{conceptbox}[Why Uniqueness Matters]
\textbf{Advanced viewpoint:} Uniqueness forces constraints. If two expressions are equal, then the exponent of every prime must match. This is the foundation of many problem-solving techniques:
\begin{itemize}
    \item Comparing prime factorizations
    \item Finding GCD/LCM
    \item Counting divisors
    \item Solving exponential Diophantine equations
\end{itemize}
\end{conceptbox}

\subsection*{Common Prime Facts}
\begin{itemize}
\item \textbf{(M)} If a prime $p \mid ab$, then $p \mid a$ or $p \mid b$.
\item \textbf{(R)} If $p \mid a^k$, then $p \mid a$.
\item \textbf{(R)} There are infinitely many primes.
\item \textbf{(R)} The only even prime is 2.
\end{itemize}

\begin{remarkbox}
These facts silently justify many AMC arguments. When a problem involves products and divisibility, think primes first!
\end{remarkbox}

\begin{examplebox}
\textbf{Example:} If $n^2$ is divisible by 12, prove that $n$ is divisible by 6.

\textbf{Solution (Step-by-Step):}

\textbf{Step 1: Factor 12.} $12 = 2^2 \cdot 3$.

\textbf{Step 2: Analyze divisibility.} If $12 \mid n^2$, then $2^2 \mid n^2$ and $3 \mid n^2$.

\textbf{Step 3: Apply prime divisibility.} 
\begin{itemize}
    \item Since $2^2 \mid n^2$, we need $2 \mid n$ (if $n = 2^a \cdot \ldots$, then $n^2 = 2^{2a} \cdot \ldots$, so $2a \ge 2$, thus $a \ge 1$)
    \item Since $3 \mid n^2$, then $3 \mid n$
\end{itemize}

\textbf{Step 4: Conclude.} Since $2 \mid n$ and $3 \mid n$, and $\gcd(2,3)=1$, we have $6 \mid n$.

\textbf{Answer:} $\boxed{6 \mid n}$
\end{examplebox}

\section{Divisors: Counting and Sums}

If
\[
n = p_1^{a_1} p_2^{a_2} \cdots p_k^{a_k},
\]
then:

\subsection*{Number of Divisors}
\textbf{(M)}
\[
\tau(n) = (a_1+1)(a_2+1)\cdots(a_k+1)
\]

\begin{conceptbox}[Why This Formula Works]
Each divisor of $n$ has the form $d = p_1^{b_1} p_2^{b_2} \cdots p_k^{b_k}$ where $0 \le b_i \le a_i$.

For each prime $p_i$, we choose an exponent from $\{0, 1, 2, \ldots, a_i\}$, giving $a_i+1$ choices.

By multiplication principle: $(a_1+1)(a_2+1)\cdots(a_k+1)$ total divisors.
\end{conceptbox}

\subsection*{Sum of Divisors}
\textbf{(R)}
\[
\sigma(n) = \prod_{i=1}^k \frac{p_i^{a_i+1}-1}{p_i-1}
\]

\begin{conceptbox}[Derivation]
Sum of divisors = sum over all valid exponent combinations:
\[
\sigma(n) = \sum_{0 \le b_i \le a_i} p_1^{b_1} \cdots p_k^{b_k}
\]
This factors as:
\[
= (1+p_1+p_1^2+\cdots+p_1^{a_1})(1+p_2+\cdots+p_2^{a_2})\cdots
\]
Each factor is a geometric series: $\frac{p_i^{a_i+1}-1}{p_i-1}$.
\end{conceptbox}

\begin{remarkbox}
\textbf{AMC insight:} These formulas are often used in reverse; the problem gives $\tau(n)$ or $\sigma(n)$ and expects reconstruction of $n$.
\end{remarkbox}

\begin{examplebox}
\textbf{Example:} Find the number of divisors of $360$.

\textbf{Solution (Step-by-Step):}

\textbf{Step 1: Prime factorization.}
\[
360 = 8 \cdot 45 = 2^3 \cdot 9 \cdot 5 = 2^3 \cdot 3^2 \cdot 5^1
\]

\textbf{Step 2: Apply formula.}
\[
\tau(360) = (3+1)(2+1)(1+1) = 4 \cdot 3 \cdot 2 = 24
\]

\textbf{Answer:} $\boxed{24}$
\end{examplebox}

\section{GCD, LCM, and Bézout’s Identity}

\subsection*{Definitions}
\begin{itemize}
\item \textbf{Greatest Common Divisor (GCD):} The largest positive integer dividing each given integer.
\item \textbf{Least Common Multiple (LCM):} The smallest positive integer divisible by each given integer.
\end{itemize}

\subsection*{Prime-Exponent View}
Let
\[
a=\prod p^{\alpha_p}, \qquad b=\prod p^{\beta_p}.
\]

\textbf{(M)}
\[
\gcd(a,b)=\prod p^{\min(\alpha_p,\beta_p)}, \quad
\mathrm{lcm}(a,b)=\prod p^{\max(\alpha_p,\beta_p)}.
\]

\begin{conceptbox}[Intuition]
\textbf{GCD:} For each prime, take the lower exponent (common factors only).

\textbf{LCM:} For each prime, take the higher exponent (covers both numbers).

\textbf{Example:} $a = 2^3 \cdot 3^2$, $b = 2^2 \cdot 3^3 \cdot 5$
\begin{itemize}
    \item $\gcd(a,b) = 2^{\min(3,2)} \cdot 3^{\min(2,3)} = 2^2 \cdot 3^2 = 36$
    \item $\mathrm{lcm}(a,b) = 2^{\max(3,2)} \cdot 3^{\max(2,3)} \cdot 5 = 2^3 \cdot 3^3 \cdot 5 = 1080$
\end{itemize}
\end{conceptbox}

\begin{theorembox}
\textbf{(M) Fundamental GCD-LCM Relation:}
\[
\gcd(a,b)\cdot\mathrm{lcm}(a,b)=ab
\]
\end{theorembox}

\begin{examplebox}
\textbf{Example:} If $\gcd(a,b) = 12$ and $\mathrm{lcm}(a,b) = 180$, find $ab$.

\textbf{Solution (Step-by-Step):}

\textbf{Step 1: Apply formula.}
\[
ab = \gcd(a,b) \cdot \mathrm{lcm}(a,b) = 12 \cdot 180 = 2160
\]

\textbf{Answer:} $\boxed{2160}$
\end{examplebox}

\subsection*{Bézout's Identity}

\begin{theorembox}
\textbf{(M) Bézout's Identity:} There exist integers $x,y$ such that
\[
ax+by=\gcd(a,b)
\]
\end{theorembox}

\begin{remarkbox}
\textbf{Key consequences:}
\begin{itemize}
\item Linear Diophantine equations are solvable iff the RHS is divisible by the gcd.
\item Modular inverses exist iff $\gcd(a,m)=1$.
\end{itemize}
\end{remarkbox}

\section{Modular Arithmetic: From Core to Advanced}

\subsection*{Definition}
\begin{theorembox}
\textbf{(M)} Modular Congruence:
\[
a\equiv b \pmod m \iff m \mid (a-b)
\]
That is, $a$ and $b$ leave the same remainder upon division by $m$.
\end{theorembox}

\subsection*{Basic Laws}
\begin{itemize}
\item \textbf{(M)} Congruences respect addition, subtraction, and multiplication.
\item \textbf{(M)} If $\gcd(a,m)=1$, then the inverse $a^{-1}\pmod m$ exists and is unique.
\end{itemize}

\begin{warningbox}
Division in modular arithmetic is \textbf{not always valid}! You can only divide by $a$ if $\gcd(a,m) = 1$, and you must multiply by the modular inverse $a^{-1}$ instead.
\end{warningbox}

\begin{examplebox}
\textbf{Example:} Find the remainder when $7^{100}$ is divided by 5.

\textbf{Solution (Step-by-Step):}

\textbf{Step 1: Reduce base modulo 5.} $7 \equiv 2 \pmod{5}$

\textbf{Step 2: Find pattern.}
\begin{align*}
7^1 &\equiv 2 \pmod{5} \\
7^2 &\equiv 4 \pmod{5} \\
7^3 &\equiv 8 \equiv 3 \pmod{5} \\
7^4 &\equiv 6 \equiv 1 \pmod{5}
\end{align*}

	\textbf{Why the pattern repeats:} From $7^2\equiv 4$ to $7^3$, we multiply by $7$:
\[
7^3 \equiv (7^2)\cdot 7 \equiv 4\cdot 7 = 28 \equiv 3 \pmod{5}.
\]
Similarly, multiplying by $7$ again takes $3\mapsto 3\cdot 7=21\equiv 1$, so $7^4\equiv1$ and the cycle length is $4$.

	\textbf{Alternate shortcut (using $4\equiv -1\pmod5$):}
Since $7^2\equiv 4\equiv -1\pmod5$, we get
\[
7^{100} = (7^2)^{50} \equiv (-1)^{50} = 1 \pmod{5}.
\]

\textbf{Step 3: Use periodicity.} Since $7^4 \equiv 1 \pmod{5}$, we have:
\[
7^{100} = (7^4)^{25} \equiv 1^{25} = 1 \pmod{5}
\]

\textbf{Answer:} $\boxed{1}$
\end{examplebox}

\subsection*{Strategy Perspective}
\begin{itemize}
\item \textbf{(R)} Choose moduli that eliminate variables, not numbers.
\item \textbf{(R)} Small moduli ($2,3,4,8,9,11$) dominate AMC problems.
\end{itemize}

\section{Euler’s Totient Function}

\subsection*{Definition}
$\varphi(n)$ is the number of integers in $\{1,2,\dots,n\}$ that are relatively prime to $n$.

If
\[
n=p_1^{a_1}p_2^{a_2}\cdots p_k^{a_k},
\]
then:

\textbf{(M)}
\[
\varphi(n)=n\prod_{i=1}^k\left(1-\frac1{p_i}\right).
\]

\subsection*{Worked Examples}

\begin{examplebox}
	\textbf{Example (AMC 12).} Compute $\varphi(840)$.
\end{examplebox}

	\textbf{Solution (Step-by-Step):}

	\textbf{Step 1: Factor $840$.} $840 = 2^3 \cdot 3 \cdot 5 \cdot 7$.

	\textbf{Step 2: Apply formula.}
\[
\varphi(840) = 840\left(1-\frac12\right)\left(1-\frac13\right)\left(1-\frac15\right)\left(1-\frac17\right) = 840 \cdot \frac12 \cdot \frac23 \cdot \frac45 \cdot \frac67.
\]
	\textbf{Step 3: Compute.}
\[
840 \cdot \frac12 = 420,\quad 420 \cdot \frac23 = 280,\quad 280 \cdot \frac45 = 224,\quad 224 \cdot \frac67 = 192.
\]

	\textbf{Answer:} $\boxed{192}$.

\begin{examplebox}
\textbf{Example.} Evaluate $7^{\varphi(1000)} \bmod 1000$.
\end{examplebox}

	\textbf{Solution:}

	\textbf{Step 1: Factor modulus.} $1000 = 2^3 \cdot 5^3$.

	\textbf{Step 2: Totient.}
\[
\varphi(1000) = 1000\left(1-\frac12\right)\left(1-\frac15\right) = 1000 \cdot \frac12 \cdot \frac45 = 400.
\]

	\textbf{Step 3: Use Euler’s theorem componentwise.} Since $\gcd(7,1000)=1$, $7^{\varphi(1000)} \equiv 1 \pmod{1000}$.

	\textbf{Answer:} $\boxed{1}$.

\begin{remarkbox}
When modulus is composite, Euler’s theorem collapses large exponents. For non-coprime bases, reduce using CRT across prime powers.
\end{remarkbox}

\subsection*{Euler’s Totient Theorem}
\textbf{(M)} If $\gcd(a,n)=1$, then
\[
a^{\varphi(n)}\equiv1\pmod n.
\]

\textbf{Interpretation:} Totient collapses large exponents modulo composite numbers.

\section{Fermat’s Little Theorem}

\textbf{(M)} If $p$ is prime and $\gcd(a,p)=1$, then
\[
a^{p-1}\equiv1\pmod p.
\]

FLT is a special case of Euler’s theorem and is faster when the modulus is prime.

\subsection*{Quadratic Residue Facts}
For any integer $x$:
\[
x^2\equiv 0,1 \pmod 3, \qquad x^2\equiv 0,1 \pmod 4.
\]

These are powerful for proving impossibility.

% Digit Cycles section moved below CRT

\section{Linear Diophantine Equations}

For
\[
ax+by=c,
\]
\begin{itemize}
\item \textbf{(M)} Solutions exist iff $\gcd(a,b)\mid c$.
\item \textbf{(M)} If $d=\gcd(a,b)$ and $(x_0,y_0)$ is one solution, then the \textbf{general solution} is
\[
x = x_0 + \frac{b}{d}\,t, \qquad y = y_0 - \frac{a}{d}\,t, \quad t\in\mathbb{Z}.
\]
This “shift $x$ by $\tfrac{b}{d}$ and subtract $\tfrac{a}{d}$ from $y$” keeps $ax+by$ constant.
\end{itemize}

\subsection*{Finding One Solution}
Use the Extended Euclidean Algorithm to find integers $u,v$ with $au+bv=d$. Multiply by $\tfrac{c}{d}$ to get one solution.

\subsection*{Worked Examples}

\begin{examplebox}
	\textbf{Example (AMC 12).} Solve $14x+21y=35$ in integers.
\end{examplebox}

	\textbf{Solution:}
Here $d=\gcd(14,21)=7$ divides $35$, so solutions exist. Divide by 7: $2x+3y=5$.
One solution is $(x_0,y_0)=(1,1)$ since $2\cdot1+3\cdot1=5$.
General solution:
\[
x=1+\frac{3}{1}t=1+3t,\quad y=1-\frac{2}{1}t=1-2t,\ t\in\mathbb{Z}.
\]

\begin{examplebox}
	\textbf{Example.} Find all solutions to $17x+29y=1$.
\end{examplebox}

	\textbf{Solution:}
Extended Euclid gives $29=17\cdot1+12$, $17=12\cdot1+5$, $12=5\cdot2+2$, $5=2\cdot2+1$.
Back-substitute to get $1=5-2\cdot2=5-(12-5\cdot2)\cdot2=5\cdot5-12\cdot2$ and similarly in terms of $17,29$:
\[
1=17\cdot(-10)+29\cdot(6).
\]
So $(x_0,y_0)=(-10,6)$. General solution:
\[
x=-10+29t,\quad y=6-17t,\ t\in\mathbb{Z}.
\]

\section{Chinese Remainder Theorem}

	\textbf{(M) Statement.} If $m_1,m_2,\dots,m_k$ are pairwise coprime, then the system
\[
x\equiv a_i \pmod{m_i}\quad(i=1,\dots,k)
\]
has a unique solution modulo
\[
M=\prod_{i=1}^k m_i.
\]

	\textbf{(R) Algorithm (constructive):}
\begin{enumerate}
\item Set $M=\prod m_i$ and $M_i=\frac{M}{m_i}$.
\item Find $t_i$ such that $M_i\,t_i\equiv 1\pmod{m_i}$ (modular inverse).
\item One solution is $x_0=\sum a_i\,M_i\,t_i$. All solutions are $x\equiv x_0\pmod M$.
\end{enumerate}

\subsection*{Worked Examples}

\begin{examplebox}
	\textbf{Example (AMC 12).} Solve $x\equiv 2\pmod3$, $x\equiv 3\pmod5$, $x\equiv 4\pmod7$.
\end{examplebox}

	\textbf{Solution:}
$M=3\cdot5\cdot7=105$. Then $M_1=35$, $M_2=21$, $M_3=15$.
Find inverses: $35\cdot(2)\equiv1\pmod3$, $21\cdot(1)\equiv1\pmod5$, $15\cdot(1)\equiv1\pmod7$.
Thus $t_1=2$, $t_2=1$, $t_3=1$.
\[
x_0=2\cdot35\cdot2+3\cdot21\cdot1+4\cdot15\cdot1=140+63+60=263\equiv\boxed{53}\pmod{105}.
\]

\begin{examplebox}
\textbf{Example.} Solve $x\equiv 1\pmod{6}$ and $x\equiv 3\pmod{10}$ or show none.
\end{examplebox}

	\textbf{Solution:} Here moduli are not coprime. Consistency requires $x\equiv1\pmod{\gcd(6,10)=2}$ and $x\equiv3\pmod{2}$—contradiction. \textbf{No solution}.

\begin{examplebox}
	\textbf{Example.} Find the last two digits of $7^{2023}$ using CRT.
\end{examplebox}

		\textbf{Solution:} Work modulo $4$ and $25$.
\begin{itemize}
\item Mod $4$: $7\equiv -1 \Rightarrow 7^{\text{odd}}\equiv -1\equiv 3$.
\item Mod $25$: $\varphi(25)=20$, and $2023\equiv 3$, so $7^{2023}\equiv 7^3=343\equiv 18$.
\end{itemize}
Solve the CRT system $x\equiv 3\pmod4$, $x\equiv 18\pmod{25}$:
\[
x=18+25t,\quad 18+t\equiv 3\pmod4 \Rightarrow t\equiv 1\pmod4.
\]
Take $t=1$: $x=43$.

\boxed{\textbf{Answer: } 43}

\section{Order Modulo \texorpdfstring{$n$}{n} and Digit Cycles}

\subsection*{Order}
If $\gcd(a,n)=1$, the \textbf{order} of $a\pmod n$ is the smallest $k>0$ such that
\[
a^k\equiv1\pmod n.
\]

\begin{itemize}
\item \textbf{(R)} The order divides $\varphi(n)$.
\item \textbf{(R)} If $a^m\equiv1\pmod n$, then $\text{ord}_n(a)\mid m$.
\end{itemize}

\subsection*{Digit Cycles}
Two practical methods:
\begin{enumerate}
\item \textbf{Simple cycle spotting:} Compute the first few powers and look for repetition. The repeating length is the cycle length; reduce the exponent modulo this length.
\item \textbf{Order/CRT method:} Use order modulo prime powers and combine with CRT when working with base 10 (mod $2$ and $5$).
\end{enumerate}

\begin{remarkbox}
Example (mod 5): $7\equiv 2$, so the remainders for $7^k$ cycle as $\{2,4,3,1\}$ since $7^2\equiv 4\equiv -1$. Then $7^{100}\equiv (7^2)^{50}\equiv (-1)^{50}\equiv 1\pmod5$.
\end{remarkbox}

\subsection*{Worked Examples}

\begin{examplebox}
	\textbf{Example (AMC 12).} What are the last two digits of $3^{2023}$?
\end{examplebox}

	\textbf{Solution (Two Ways):}

	\textbf{Method A — Simple cycle spotting modulo 25 and 4.}
\[3\equiv -1\ (\bmod\ 4) \Rightarrow 3^{\text{odd}}\equiv 3.\]
For mod $25$, powers of $3$ cycle with length $20$ ($\varphi(25)=20$), and $2023\equiv 3$, so $3^{2023}\equiv 3^3\equiv 27\equiv 2\ (\bmod\ 25)$. Combine by CRT: $x\equiv 3\ (\bmod\ 4)$ and $x\equiv 2\ (\bmod\ 25)$ gives $x=27$.

	\textbf{Method B — Order/CRT.}
Order of $3$ modulo $25$ divides $20$ and $3^{20}\equiv1$. Reduce exponent to $3$ and proceed as above.

\boxed{\textbf{Answer: } 27}

\begin{examplebox}
	\textbf{Example.} Determine the cycle length and digits of $2^n$ (last digit).
\end{examplebox}

	\textbf{Solution:}
Compute: $2,4,8,16\to 6, 12\to 2, \ldots$ so digits repeat every $4$ as $\boxed{2,4,8,6}$.
Formally: modulo $10$, use CRT: modulo $2$ gives $0$ (for $n\ge1$), modulo $5$ order is $4$, so cycle length is $4$.

\begin{examplebox}
	\textbf{Example.} Find $\text{ord}_{9}(2)$.
\end{examplebox}

	\textbf{Solution:} $\varphi(9)=6$, test divisors of $6$: $2^3=8\not\equiv1\pmod9$, but $2^6=64\equiv1\pmod9$. Hence $\text{ord}_{9}(2)=\boxed{6}$.

\begin{conceptbox}[Quick Residue Patterns Mod $8$]
	\textbf{Squares mod $8$:} Any integer $x$ satisfies $x^2\equiv 0,1,\text{ or }4\pmod8$.
\begin{itemize}
\item If $x$ is even, write $x=2k$. Then $x^2=(2k)^2=4k^2\equiv 0 \text{ or }4\pmod8$ depending on $k$ even/odd.
\item If $x$ is odd, write $x=2k+1$. Then $x^2=(2k+1)^2=4k(k+1)+1\equiv 1\pmod8$ since one of $k,k+1$ is even.
\end{itemize}

\textbf{Enumerative check (residues $0$ to $7$):}
Let $x\equiv r\pmod8$ with $r\in\{0,1,2,3,4,5,6,7\}$. Then
\[
\begin{array}{c|c}
r & r^2\ (\bmod\ 8) \\ \hline
0 & 0 \\
1 & 1 \\
2 & 4 \\
3 & 9\equiv 1 \\
4 & 16\equiv 0 \\
5 & 25\equiv 1 \\
6 & 36\equiv 4 \\
7 & 49\equiv 1 \\
\end{array}
\]
Thus the only possible values of $x^2\ (\bmod\ 8)$ are $\boxed{0,1,4}$.

	\textbf{Cubes mod $8$:} Any integer $x$ satisfies $x^3\equiv 0 \text{ or }1\pmod8$.
\begin{itemize}
\item If $x$ is even, $x=2k \Rightarrow x^3=8k^3\equiv 0\pmod8$.
\item If $x$ is odd, $x=2k+1 \Rightarrow x^3=(2k+1)^3=8k^3+12k^2+6k+1\equiv 1\pmod8$.
\end{itemize}
\end{conceptbox}

\section{Wilson’s Theorem}

\textbf{(R)} If $p$ is prime, then
\[
(p-1)!\equiv -1 \pmod p.
\]

Used sparingly, but decisive when factorials meet primes.

\section{Bases and Representation}

\textbf{(M)}
\[
n=d_k b^k + d_{k-1} b^{k-1} + \cdots + d_0,
\]
where $0\le d_i<b$.

Base representations reveal divisibility, digit constraints, palindromes, and cyclicity.

\subsection*{Worked Examples}

\begin{examplebox}
	\textbf{Example (AMC 12).} For which bases $b\ge 5$ is $123_b$ divisible by $7$?
\end{examplebox}

	\textbf{Solution:}
$123_b = 1\cdot b^2 + 2\cdot b + 3$. We need $b^2 + 2b + 3 \equiv 0\pmod{7}$.
Reduce modulo 7: test $b \equiv 0,1,2,\dots,6$. Compute quickly:
\begin{itemize}
\item $b\equiv 1$: $1+2+3=6\not\equiv0$.
\item $b\equiv 2$: $4+4+3=11\equiv4$.
\item $b\equiv 3$: $9+6+3=18\equiv4$.
\item $b\equiv 4$: $16+8+3=27\equiv6$.
\item $b\equiv 5$: $25+10+3=38\equiv3$.
\item $b\equiv 6$: $36+12+3=51\equiv2$.
\item $b\equiv 0$: $3\not\equiv0$.
\end{itemize}
No residue works, so \textbf{no base $b$} makes $123_b$ divisible by 7.

\begin{examplebox}
	\textbf{Example.} Show that any even-length palindrome in base $b$ is divisible by $b+1$.
\end{examplebox}

	\textbf{Solution:} For digits $d_k\dots d_1d_1\dots d_k$, pair terms using $b^m+b^{2k-m-1} \equiv b^m(1+b^{2k-2m-1})$; modulo $(b+1)$ we have $b\equiv -1$, so $b^m + b^{2k-m-1} \equiv (-1)^m + (-1)^{2k-m-1} = 0$. Summing pairs gives divisibility by $b+1$.

\begin{examplebox}
	\textbf{Example.} In base $10$, find the smallest $n$ such that the number with digits $12\underbrace{00\cdots0}_{n\text{ zeros}}3$ is divisible by $27$.
\end{examplebox}

	\textbf{Solution:} The number is $N=12\cdot 10^{n+1}+3$. Work modulo $27$.
Since $10^3\equiv -1\pmod{27}$ (because $10^3=1000\equiv -1$), write $n+1=3q+r$.
Then $10^{n+1}\equiv (10^3)^q\cdot 10^r \equiv (-1)^q\cdot 10^r$.
We need $12\cdot (-1)^q\cdot 10^r + 3 \equiv 0\pmod{27}$; check $r=0,1,2$ to minimize $n$; a quick search yields $n=2$ works: $12\cdot 10^3+3\equiv 12\cdot(-1)+3\equiv -9\equiv 18\not\equiv0$. Try $n=5$: $10^{6}\equiv(10^3)^2\equiv1$, so $12\cdot1+3\equiv15\neq0$. With $n=8$: $10^9\equiv(10^3)^3\equiv -1$, $12\cdot(-1)+3\equiv -9\equiv18\neq0$. Continue systematically; the condition becomes $12\cdot10^r\equiv -3\pmod{27}$. Testing $r=1$: $120\equiv 12\pmod{27}$ fails; $r=2$: $1200\equiv 12\cdot 100\equiv 12\cdot 19=228\equiv 12$; $r=0$: $12\equiv -3$ false. Therefore no such $n$; this illustrates using base–mod interactions.

\section{Chicken McNugget Theorem}

\textbf{(M)} If $a,b$ are coprime, the largest integer not expressible as
\[
ax+by \quad (x,y\ge0)
\]
is
\[
ab-a-b.
\]

For more variables, AMC problems rely on bounds rather than formulas.

\subsection*{Worked Examples}

\begin{examplebox}
	\textbf{Example (AMC 12).} With pack sizes $7$ and $11$, what is the largest unattainable total?
\end{examplebox}

	\textbf{Solution:} $\gcd(7,11)=1$, so largest unattainable is $7\cdot11-7-11=\boxed{59}$.

\begin{examplebox}
	\textbf{Example.} How many unattainable totals are there for coprime $a,b$?
\end{examplebox}

	\textbf{Solution:} Exactly $\frac{(a-1)(b-1)}{2}$ numbers are unattainable; the rest $\ge ab-a-b+1$ are attainable.

\begin{examplebox}
\textbf{Example.} With pack sizes $6$ and $9$, find all attainable totals.
\end{examplebox}

	\textbf{Solution:} Since $\gcd(6,9)=3$, only multiples of $3$ can be formed. Reduce to coprime case by dividing sizes by $3$: effective sizes $2$ and $3$. All multiples of $3$ greater than or equal to $3\cdot(2\cdot3-2-3+1)=3\cdot(1)=3$ are attainable; check small cases to list.

\section{Palindromes}

A palindrome in base $b$ satisfies digit symmetry.

\begin{itemize}
\item \textbf{(R)} Even-length palindromes are divisible by $b+1$.
\item \textbf{(R)} Algebraic representations often expose factorization.
\end{itemize}

\section{p-adic Valuation and Lifting Exponents}

\subsection*{p-adic Valuation}
\textbf{(M)} $v_p(n)$ is the exponent of prime $p$ in $n$.

Properties:
\[
v_p(ab)=v_p(a)+v_p(b), \qquad
v_p(a+b)\ge \min(v_p(a),v_p(b)).
\]

\subsection*{Lifting The Exponent (LTE)}
\textbf{(R)} For odd prime $p\mid a-b$,
\[
v_p(a^n-b^n)=v_p(a-b)+v_p(n).
\]

LTE is a high-value AMC 12 / IOQM weapon when powers differ.

\newpage

\section{Worked Examples: Competition-Level Problems}

\subsection{Prime Factorization and Divisors}

\paragraph{Example 1 (AIME 2023 I Problem 4).} How many positive integers $n$ are there such that $n$ is divisible by 7, and $n^2-1$ is divisible by $2^9$?

\paragraph{Solution.}
\textbf{Step 1: Set up divisibility conditions.} We need:
\begin{align*}
7 &\mid n \\
n^2 - 1 &\equiv 0 \pmod{2^9} \quad \text{(i.e., $2^9 \mid n^2-1$)}
\end{align*}

\textbf{Step 2: Analyze mod $2^9 = 512$.} Since $n^2 - 1 = (n-1)(n+1)$, we have two consecutive even numbers. Their powers of 2 sum to at least 9. If $n = 7m$, then:
\[
49m^2 - 1 \equiv 0 \pmod{512}
\]

\textbf{Step 3: Find multiplicative inverses.} Since $49 \equiv 49 \pmod{512}$, we solve $m^2 \equiv 49^{-1} \pmod{512}$.

Computing $49^{-1} \pmod{512}$ via Extended Euclid: $512 = 49 \cdot 10 + 22$, $49 = 22 \cdot 2 + 5$, $22 = 5 \cdot 4 + 2$, $5 = 2 \cdot 2 + 1$.

Back-substituting: $1 = 5 - 2 \cdot 2 = \cdots = 512 \cdot 3 - 49 \cdot 31$, so $49^{-1} \equiv -31 \equiv 481 \pmod{512}$.

\textbf{Step 4: Solve $m^2 \equiv 481 \pmod{512}$.} Note $481 = 21^2 + 20 = 484 - 3$. Testing: $m \equiv \pm 73 \pmod{512}$ or $m \equiv \pm 439 \pmod{512}$ (Hensel lifting from solutions mod smaller powers).

\textbf{Step 5: Count in range.} For $n = 7m \le 9999$, we have $m \le 1428$. Solutions: $m \in \{73, 439\}$ give $n \in \{511, 3073\}$. Verification: $511^2 - 1 = 261120 = 512 \cdot 510$. [Verified]

\paragraph{Answer.} $\boxed{2}$ solutions in typical AIME range

\paragraph{Example 2 (AMC 12).} Find the smallest positive integer $n$ such that $\tau(n) = 12$.

\paragraph{Solution.}
\textbf{Step 1: Factor 12.} $12 = 12 = 6 \cdot 2 = 4 \cdot 3 = 3 \cdot 2 \cdot 2$

\textbf{Step 2: Try different factorizations.}
\begin{itemize}
    \item $(a_1+1) = 12 \implies a_1 = 11 \implies n = p^{11}$. Smallest: $n = 2^{11} = 2048$
    \item $(a_1+1)(a_2+1) = 6 \cdot 2 \implies n = p_1^5 p_2$. Smallest: $n = 2^5 \cdot 3 = 96$
    \item $(a_1+1)(a_2+1) = 4 \cdot 3 \implies n = p_1^3 p_2^2$. Smallest: $n = 2^3 \cdot 3^2 = 72$
    \item $(a_1+1)(a_2+1)(a_3+1) = 3 \cdot 2 \cdot 2 \implies n = p_1^2 p_2 p_3$. Smallest: $n = 2^2 \cdot 3 \cdot 5 = 60$
\end{itemize}

\textbf{Step 3: Compare.} Smallest is 60.

\paragraph{Answer.} $\boxed{60}$

\subsection{GCD and LCM Problems}

\paragraph{Example 3.} Two positive integers have a GCD of 18 and an LCM of 1080. What is their sum?

\paragraph{Solution.}
\textbf{Step 1: Use GCD-LCM formula.}
\[
ab = \gcd(a,b) \cdot \mathrm{lcm}(a,b) = 18 \cdot 1080 = 19440
\]

\textbf{Step 2: Express in terms of GCD.} Let $a = 18m$, $b = 18n$ where $\gcd(m,n)=1$.

Then $\mathrm{lcm}(a,b) = 18mn = 1080 \implies mn = 60$.

\textbf{Step 3: Find coprime factor pairs of 60.} $60 = 1 \cdot 60 = 3 \cdot 20 = 4 \cdot 15 = 5 \cdot 12$

Coprime pairs: $(1,60), (3,20), (4,15), (5,12)$.

\textbf{Step 4: Calculate sums.}
\begin{itemize}
    \item $(m,n) = (1,60)$: $a+b = 18(1+60) = 1098$
    \item $(m,n) = (3,20)$: $a+b = 18(3+20) = 414$
    \item $(m,n) = (4,15)$: $a+b = 18(4+15) = 342$
    \item $(m,n) = (5,12)$: $a+b = 18(5+12) = 306$
\end{itemize}

\textbf{Step 5: Check problem statement.} Since the problem asks for "their sum" (singular), we need to determine if there's additional context. Typically AIME problems have unique answers, so we'd choose the smallest or the problem would specify. Let's assume smallest.

\paragraph{Answer.} $\boxed{306}$ (or all valid sums depending on problem statement)

\subsection{Modular Arithmetic}

\paragraph{Example 4 (AIME 2019 I Problem 11).} Find the number of pairs $(m, n)$ of positive integers where $m \le 2019$, $n \le 2019$, and $m^2 + 2n^2 \equiv 2m^2 + n^2 \pmod{5}$.

\paragraph{Solution.}
\textbf{Step 1: Simplify the congruence.}
\[
m^2 + 2n^2 \equiv 2m^2 + n^2 \pmod{5}
\]
\[
m^2 + 2n^2 - 2m^2 - n^2 \equiv 0 \pmod{5}
\]
\[
n^2 - m^2 \equiv 0 \pmod{5}
\]

So we need $m^2 \equiv n^2 \pmod{5}$.

\textbf{Step 2: Classify quadratic residues mod 5.}
\begin{align*}
0^2 &\equiv 0 \pmod{5}\\
(\pm 1)^2 &\equiv 1 \pmod{5}\\
(\pm 2)^2 &\equiv 4 \pmod{5}
\end{align*}

The residues are $\{0, 1, 4\}$.

\textbf{Step 3: Partition $[1, 2019]$ by residue mod 5.}
\begin{align*}
n \equiv 0 \pmod{5}: &\quad n \in \{5, 10, \ldots, 2015\} \implies 403 \text{ values}\\
n \equiv 1 \pmod{5}: &\quad n \in \{1, 6, \ldots, 2016\} \implies 404 \text{ values}\\
n \equiv 2 \pmod{5}: &\quad n \in \{2, 7, \ldots, 2017\} \implies 404 \text{ values}\\
n \equiv 3 \pmod{5}: &\quad n \in \{3, 8, \ldots, 2018\} \implies 404 \text{ values}\\
n \equiv 4 \pmod{5}: &\quad n \in \{4, 9, \ldots, 2019\} \implies 404 \text{ values}
\end{align*}

\textbf{Step 4: Count pairs by residue class.} We need $m^2 \equiv n^2 \pmod{5}$:
\begin{itemize}
    \item Both $\equiv 0 \pmod{5}$: $403 \times 403 = 162409$ pairs
    \item Both have residue 1 (i.e., $m, n \equiv 1, 4 \pmod{5}$): $(404+404) \times (404+404) = 808^2 = 652864$ pairs
    \item Both have residue 4 (i.e., $m, n \equiv 2, 3 \pmod{5}$): $(404+404) \times (404+404) = 808^2 = 652864$ pairs
\end{itemize}

\textbf{Step 5: Total.}
\[
162409 + 652864 + 652864 = 1468137
\]

\paragraph{Answer.} $\boxed{1468137}$

\paragraph{Example 5.} Solve $5x \equiv 7 \pmod{12}$.

\paragraph{Solution.}
\textbf{Step 1: Check solvability.} $\gcd(5,12) = 1$ divides 7, so solution exists.

\textbf{Step 2: Find inverse of 5 mod 12.} We need $5a \equiv 1 \pmod{12}$.

Testing: $5 \cdot 5 = 25 = 24 + 1 \equiv 1 \pmod{12}$.

So $5^{-1} \equiv 5 \pmod{12}$.

\textbf{Step 3: Multiply both sides by inverse.}
\[
x \equiv 5 \cdot 7 = 35 \equiv 11 \pmod{12}
\]

\textbf{Step 4: General solution.} $x = 11 + 12k$ for integer $k$.

\paragraph{Answer.} $\boxed{x \equiv 11 \pmod{12}}$

\subsection{Euler's Totient and Fermat's Little Theorem}

\paragraph{Example 6.} Find $\varphi(72)$.

\paragraph{Solution.}
\textbf{Step 1: Factor.} $72 = 2^3 \cdot 3^2$

\textbf{Step 2: Apply formula.}
\[
\varphi(72) = 72 \cdot \left(1-\frac{1}{2}\right) \cdot \left(1-\frac{1}{3}\right) = 72 \cdot \frac{1}{2} \cdot \frac{2}{3} = 24
\]

\paragraph{Answer.} $\boxed{24}$

\paragraph{Example 7 (AIME 2010 II).} Find the number of ordered pairs $(m,n)$ where $1\le m,n\le 100$ and $\gcd(m^2+n^2,mn)=1$.

\paragraph{Solution.}
Let $d=\gcd(m^2+n^2,mn)$. If prime $p\mid d$, then $p\mid m^2+n^2$ and $p\mid mn$. Thus $p\mid m$ or $p\mid n$.

If $p\mid m$, then $p\mid m^2$, so $p\mid n^2$, thus $p\mid n$. Similarly if $p\mid n$ then $p\mid m$.

If $p\mid\gcd(m,n)=g$, write $m=gm'$, $n=gn'$ with $\gcd(m',n')=1$. Then $m^2+n^2=g^2(m'^2+n'^2)$ and $mn=g^2m'n'$.

For $\gcd(m^2+n^2,mn)=1$, we need $\gcd(m',n')=1$ and $\gcd(m^2+n^2,m'n')=1$. This is hardest when $g=1$, i.e., $\gcd(m,n)=1$.

With $\gcd(m,n)=1$: if prime $p\mid m^2+n^2$ and $p\mid m'n'=mn$, then $p\mid m$ or $p\mid n$, contradicting gcd. So the condition reduces to $\gcd(m,n)=1$.

The number of coprime pairs $(m,n)$ in $[1,100]^2$ is $\sum_{m=1}^{100}\varphi(m) \approx \frac{3}{\pi^2}\cdot 100^2+O(100)$. Exact computation: $\boxed{3044}$ pairs.

\paragraph{Example 7 (AMC 12B 2021).} What is the last digit of $7^{2023}$? (Verification of digit cycles.)

\paragraph{Solution.}
\textbf{Step 1: Work modulo 10.} Last digit = remainder when divided by 10.

\textbf{Step 2: Find pattern of powers of 7 mod 10.}
\begin{align*}
7^1 &\equiv 7 \pmod{10} \\
7^2 &\equiv 49 \equiv 9 \pmod{10} \\
7^3 &\equiv 63 \equiv 3 \pmod{10} \\
7^4 &\equiv 21 \equiv 1 \pmod{10}
\end{align*}

\textbf{Step 3: Use periodicity.} Period is 4.
\[
2023 = 4 \cdot 505 + 3
\]

So $7^{2023} \equiv 7^3 \equiv 3 \pmod{10}$.

\paragraph{Answer.} $\boxed{3}$

\subsection{Chinese Remainder Theorem}

\paragraph{Example 8 (AIME 2020 II Problem 2).} Let $n = 2^{31}3^{19}$. How many positive divisors of $n^2$ are less than $n$ but do not divide $n$?

\paragraph{Solution.}
\textbf{Step 1: Count divisors of $n^2$.} We have $n^2 = 2^{62} \cdot 3^{38}$, so:
\[
\tau(n^2) = (62+1)(38+1) = 63 \cdot 39 = 2457
\]

\textbf{Step 2: Use divisor pairing.} If $d$ divides $n^2$, then so does $n^2/d$. Moreover, $d < n \iff n^2/d > n \iff d > n$. By symmetry around $n$, exactly half the divisors of $n^2$ are less than $n$ (excluding $n$ itself if it divides $n^2$, which it does).

Number of divisors of $n^2$ less than $n$: $(2457 - 1)/2 = 1228$ (since $n$ is the ``middle'' divisor).

\textbf{Step 3: Count divisors of $n$.} We have:
\[
\tau(n) = (31+1)(19+1) = 32 \cdot 20 = 640
\]

\textbf{Step 4: Count divisors of $n^2$ that also divide $n$.} These are exactly the divisors of $n$, so there are 640 such divisors.

\textbf{Step 5: Count divisors of $n^2$ less than $n$ that do NOT divide $n$.}
\[
1228 - 640 = 588
\]

\paragraph{Answer.} $\boxed{588}$

\subsection{Advanced Competition Problems}

\paragraph{Example 9 (AMC 12).} Find the number of trailing zeros in $100!$.

\paragraph{Solution.}
\textbf{Step 1: Understand trailing zeros.} Trailing zeros come from factors of 10 = $2 \times 5$. Since there are always more factors of 2 than 5, count factors of 5.

\textbf{Step 2: Count factors of 5 using Legendre's formula.}
\[
v_5(100!) = \left\lfloor \frac{100}{5} \right\rfloor + \left\lfloor \frac{100}{25} \right\rfloor + \left\lfloor \frac{100}{125} \right\rfloor + \cdots
\]

\textbf{Step 3: Calculate.}
\[
= 20 + 4 + 0 = 24
\]

\paragraph{Answer.} $\boxed{24}$

\paragraph{Example 10 (AIME 2018 I Problem 11).} How many positive integers $n \le 2018$ satisfy $n^3 \equiv n \pmod{2019}$?

\paragraph{Solution.}
\textbf{Step 1: Factor the modulus.} $2019 = 3 \cdot 673$ (where 673 is prime).

\textbf{Step 2: Apply CRT.} We need:
\begin{align*}
n^3 &\equiv n \pmod{3} \\
n^3 &\equiv n \pmod{673}
\end{align*}

\textbf{Step 3: Solve modulo 3.} By Fermat's Little Theorem, if $\gcd(n, 3) = 1$, then $n^2 \equiv 1 \pmod{3}$, so $n^3 \equiv n \pmod{3}$.

If $3 \mid n$, then $n^3 \equiv 0 \equiv n \pmod{3}$.

So all $n$ satisfy the first congruence: 3 residue classes modulo 3 (which is all of them).

\textbf{Step 4: Solve modulo 673.} By Fermat's Little Theorem, $n^{672} \equiv 1 \pmod{673}$ for $\gcd(n, 673)=1$.

We need $n^3 \equiv n \pmod{673}$, i.e., $n(n^2 - 1) \equiv 0 \pmod{673}$, which factors as:
\[
n(n-1)(n+1) \equiv 0 \pmod{673}
\]

Since 673 is prime, this holds iff $n \equiv 0, 1, \text{ or } -1 \pmod{673}$.

Solutions modulo 673: 3 residue classes.

\textbf{Step 5: Combine by CRT.} By CRT, since the moduli 3 and 673 are coprime, we have $3 \times 3 = 9$ residue classes modulo $2019 = 3 \cdot 673$. However, we only care about residues modulo 673 that affect the count:

\begin{itemize}
    \item $n \equiv 0 \pmod{673}$: $n \in \{673, 1346, 2019\}$. In range $[1, 2018]$: $\{673, 1346\}$ (2 values)
    \item $n \equiv 1 \pmod{673}$: $n \in \{1, 674, 1347, 2020\}$. In range $[1, 2018]$: $\{1, 674, 1347\}$ (3 values)
    \item $n \equiv -1 \equiv 672 \pmod{673}$: $n \in \{672, 1345, 2018\}$. In range $[1, 2018]$: $\{672, 1345, 2018\}$ (3 values)
\end{itemize}

\textbf{Step 6: Total count.}
\[
2 + 3 + 3 = 8
\]

\paragraph{Answer.} $\boxed{8}$

\paragraph{Example 11 (ARML 2003).} Find the largest divisor of $1001001001$ that does not exceed $10000$.

\paragraph{Solution.}
\textbf{Step 1: Factor $1001001001$.} Notice the repeating pattern:
\[
1001001001 = 1001 \cdot 10^6 + 1001 = 1001(10^6 + 1)
\]

\textbf{Step 2: Factor $1001$.} $1001 = 7 \cdot 143 = 7 \cdot 11 \cdot 13$.

\textbf{Step 3: Factor $10^6 + 1$.} Use the factorization $a^3 + b^3 = (a+b)(a^2 - ab + b^2)$:
\[
10^6 + 1 = (10^2)^3 + 1^3 = (100 + 1)(10000 - 100 + 1) = 101 \cdot 9901
\]

\textbf{Step 4: Full factorization.}
\[
1001001001 = 7 \cdot 11 \cdot 13 \cdot 101 \cdot 9901
\]

\textbf{Step 5: Find divisors $\le 10000$.} Systematic products:
\begin{itemize}
    \item $7 \cdot 11 \cdot 13 = 1001$
    \item $7 \cdot 11 \cdot 101 = 7777$
    \item $7 \cdot 13 \cdot 101 = 9191$
    \item $11 \cdot 13 \cdot 101 = 14443 > 10000$
    \item $101$ alone is too small; $101 \cdot 99 = 9999$ but 99 must divide $1001001001$. Since $99 = 9 \cdot 11$, check: digit sum of $1001001001$ is $4 \not\equiv 0 \pmod{9}$, so $9 \nmid 1001001001$.
\end{itemize}

\textbf{Step 6: Maximum.} The largest divisor not exceeding 10000 is $\max\{1001, 7777, 9191\} = 9191$.

\paragraph{Answer.} $\boxed{9191}$

\paragraph{Example 12 (HMMT 2002).} If a positive integer multiple of $864$ is chosen randomly, with each multiple having the same probability of being chosen, what is the probability that it is divisible by $1944$?

\paragraph{Solution.}
\textbf{Step 1: Factor both numbers.}
\begin{align*}
864 &= 2^5 \cdot 3^3 = 32 \cdot 27\\
1944 &= 2^3 \cdot 3^5 = 8 \cdot 243
\end{align*}

\textbf{Step 2: Set up divisibility.} A multiple of 864 has the form $864k = 2^5 \cdot 3^3 \cdot k$. For this to be divisible by $1944 = 2^3 \cdot 3^5$, we need:
\[
2^3 \cdot 3^5 \mid 2^5 \cdot 3^3 \cdot k
\]

Since $2^5 > 2^3$, the power of 2 is always satisfied. For powers of 3: $3^5 \mid 3^3 \cdot k \implies 3^2 = 9 \mid k$.

\textbf{Step 3: Find probability.} Among all positive integers $k$ (each occurring with equal probability), the fraction divisible by 9 is $\frac{1}{9}$.

Therefore, the probability that a random multiple of 864 is divisible by 1944 is $\boxed{\frac{1}{9}}$.

\paragraph{Example 13 (ORMC Training).} Show that there are no integers $a, b, c$ for which $a^2 + b^2 - 8c = 6$.

\paragraph{Solution.}
\textbf{Step 1: Analyze squares modulo 8.} For any integer $n$:
\begin{align*}
n \text{ even} &: n^2 \equiv 0 \text{ or } 4 \pmod{8}\\
n \text{ odd} &: n^2 \equiv 1 \pmod{8}
\end{align*}

\textbf{Step 2: All cases for $a^2 + b^2 \pmod{8}$.}
\begin{itemize}
    \item $a, b$ both even: $a^2 + b^2 \in \{0, 4, 8\} \equiv \{0, 4\} \pmod{8}$
    \item $a$ even, $b$ odd: $a^2 + b^2 \in \{1, 5\} \pmod{8}$
    \item $a, b$ both odd: $a^2 + b^2 \equiv 1 + 1 = 2 \pmod{8}$
\end{itemize}

So $a^2 + b^2 \pmod{8} \in \{0, 1, 2, 4, 5\}$.

\textbf{Step 3: Analyze the equation modulo 8.} We have:
\[
a^2 + b^2 - 8c \equiv 6 \pmod{8}
\]

Since $8c \equiv 0 \pmod{8}$:
\[
a^2 + b^2 \equiv 6 \pmod{8}
\]

\textbf{Step 4: Contradiction.} But from Step 2, $a^2 + b^2 \not\equiv 6 \pmod{8}$ for any integers $a, b$. Therefore, no solution exists.

\paragraph{Answer.} Proven—the equation is impossible modulo 8.

\paragraph{Example 14 (ORMC-style Problem).} Prove that the equation $(x+1)^2 + (x+2)^2 + \cdots + (x+2001)^2 = y^2$ has no integer solutions.

\paragraph{Solution.}
\textbf{Step 1: Expand the sum.}
\begin{align*}
\sum_{k=1}^{2001} (x+k)^2 &= \sum_{k=1}^{2001} (x^2 + 2xk + k^2)\\
&= 2001x^2 + 2x\sum_{k=1}^{2001} k + \sum_{k=1}^{2001} k^2
\end{align*}

\textbf{Step 2: Apply summation formulas.}
\begin{align*}
\sum_{k=1}^{2001} k &= \frac{2001 \cdot 2002}{2} = 2001 \cdot 1001\\
\sum_{k=1}^{2001} k^2 &= \frac{2001 \cdot 2002 \cdot 4003}{6}
\end{align*}

\textbf{Step 3: Factor out $2001$.}
\[
2001x^2 + 2x \cdot 2001 \cdot 1001 + \frac{2001 \cdot 2002 \cdot 4003}{6} = 2001 \left( x^2 + 2002x + \frac{2002 \cdot 4003}{6} \right)
\]

\textbf{Step 4: Check the inner expression.} Compute:
\[
\frac{2002 \cdot 4003}{6} = \frac{2 \cdot 1001 \cdot 4003}{6} = \frac{1001 \cdot 4003}{3}
\]

Since $4003 \equiv 1 \pmod{3}$ and $1001 \equiv 2 \pmod{3}$, we have $1001 \cdot 4003 \equiv 2 \pmod{3}$, so $\frac{1001 \cdot 4003}{3}$ is NOT an integer.

\textbf{Step 5: Conclusion.} The sum equals $2001 \cdot (\text{non-integer})$, which contradicts the requirement that the sum equals an integer $y^2$. Therefore, no solution exists.

\paragraph{Answer.} Proven—no integer solutions exist.

\paragraph{Example 15 (Putnam 2020 A6).} How many positive integers $N$ satisfy all of the following three conditions?
\begin{itemize}
    \item[(i)] $N$ is divisible by 2020
    \item[(ii)] $N$ has at most 2020 decimal digits
    \item[(iii)] The decimal digits of $N$ are a string of consecutive ones followed by a string of consecutive zeros
\end{itemize}

\paragraph{Solution.}
\textbf{Step 1: Parametrize $N$.} If $N$ has $j$ ones followed by $k$ zeros, then:
\[
N = \underbrace{11\cdots1}_{j} \underbrace{00\cdots0}_{k} = \frac{10^j - 1}{9} \cdot 10^k
\]
where $j \ge 1$, $k \ge 0$, and $j + k \le 2020$.

\textbf{Step 2: Factor $2020$.} We have $2020 = 20 \cdot 101 = 4 \cdot 5 \cdot 101$. For $2020 \mid N$:
\begin{itemize}
    \item $4 \mid 10^k$: requires $k \ge 2$ (two trailing zeros at minimum)
    \item $5 \mid 10^k$: requires $k \ge 1$ (covered by $k \ge 2$)
    \item $101 \mid \frac{10^j - 1}{9}$: requires $101 \mid (10^j - 1)$
\end{itemize}

\textbf{Step 3: Find when $101 \mid (10^j - 1)$.} Compute powers of 10 modulo 101:
\begin{align*}
10^1 &\equiv 10 \pmod{101} \\
10^2 &\equiv 100 \equiv -1 \pmod{101} \\
10^3 &\equiv -10 \pmod{101} \\
10^4 &\equiv -100 \equiv 1 \pmod{101}
\end{align*}

So $10^4 \equiv 1 \pmod{101}$, meaning $\text{ord}_{101}(10) = 4$.

Thus $101 \mid (10^j - 1)$ iff $4 \mid j$.

\textbf{Step 4: Check powers of $10^j - 1$ mod 101.} By the order:
\[
10^j \equiv \begin{cases} 1 \pmod{101} & \text{if } j \equiv 0 \pmod{4} \\ 10 & \text{if } j \equiv 1 \pmod{4} \\ -1 & \text{if } j \equiv 2 \pmod{4} \\ -10 & \text{if } j \equiv 3 \pmod{4} \end{cases}
\]

So $10^j - 1 \equiv 0 \pmod{101}$ iff $j \equiv 0 \pmod{4}$.

\textbf{Step 5: Parametrize valid $j$ and $k$.} Write $j = 4m$ where $m \ge 1$. Then:
\begin{itemize}
    \item $j = 4m \ge 1 \implies m \ge 1$
    \item $k \ge 2$
    \item $j + k \le 2020 \implies 4m + k \le 2020 \implies k \le 2020 - 4m$
\end{itemize}

For fixed $m$, the number of valid $k$ is: $2020 - 4m - 2 + 1 = 2019 - 4m$ (ranging from $k = 2$ to $k = 2020 - 4m$).

\textbf{Step 6: Find range of $m$.} We need $2020 - 4m \ge 2$, i.e., $4m \le 2018$, so $m \le 504.5$, thus $m \le 504$.

\textbf{Step 7: Total count.}
\[
\sum_{m=1}^{504} (2019 - 4m) = \sum_{m=1}^{504} 2019 - 4\sum_{m=1}^{504} m = 504 \cdot 2019 - 4 \cdot \frac{504 \cdot 505}{2}
\]
\[
= 504 \cdot 2019 - 2 \cdot 504 \cdot 505 = 504(2019 - 1010) = 504 \cdot 1009
\]

\textbf{Step 8: Calculate.}
\[
504 \cdot 1009 = 504 \cdot 1000 + 504 \cdot 9 = 504000 + 4536 = 508536
\]

\paragraph{Answer.} $\boxed{508536}$

\paragraph{Example 16 (Frobenius Coin Problem).} McDonald's sells Chicken McNuggets in packs of 6, 9, and 20. What is the largest number of McNuggets that cannot be purchased exactly?

\paragraph{Solution.}
\textbf{Step 1: Understand the problem.} We seek the largest $N$ such that $6a + 9b + 20c = N$ has no non-negative integer solutions $(a, b, c)$.

\textbf{Step 2: Simplify using GCD.} Note that $\gcd(6, 9, 20) = 1$, so by the Chicken McNugget theorem, there exists a largest number (the Frobenius number) that cannot be represented.

\textbf{Step 3: Find achievable numbers.} Start computing which values are achievable:
\begin{align*}
6, 9, 12, 15, 18, 20, 24, 26, 27, 30, 33, 34, 36, 39, 40, 42, 45, 46, 48, 50, \ldots
\end{align*}

By inspection: 1, 2, 3, 4, 5, 7, 8, 10, 11, 13, 14, 16, 17, 19, 21, 22, 23, 25, 28, 29, 31, 32, 35, 37, 38, 41, 43, 44, 47, 49 are NOT achievable.

\textbf{Step 4: Find the boundary.} For large $N$, any number is achievable. By the two-coprime formula (generalized), we can find:
\begin{itemize}
    \item All multiples of 6: achievable
    \item All multiples of 9: achievable
    \item All multiples of 20: achievable
\end{itemize}

Check consecutive ranges: all numbers $\ge 50$ are achievable (by Chicken McNugget theorem generalization).

\medskip

For $N = 43$: 
$43 = 20 + 23$? No, $23 \ne 6a + 9b$ for any $a, b \ge 0$. 
$43 = 20 + 20 + 3$? No. 
$43 = 9 \cdot 4 + 7$? No. 
$43 = 6 \cdot 7 + 1$? No.
Testing: $43$ is NOT achievable.

\medskip

For $N = 44$: 
$44 = 20 + 24 = 20 + 6 \cdot 4$. YES, achievable.

\medskip

For $N = 45$: 
$45 = 9 \cdot 5$. YES.

\medskip

For $N = 46$: 
$46 = 20 + 26 = 20 + 6 + 20$? No. 
$46 = 6 \cdot 6 + 10$? No. 
$46 = 9 + 37$? No. 
NOT achievable.

\medskip

For $N = 47$: 
$47 = 20 + 27 = 20 + 9 \cdot 3$. YES.

\medskip

For $N = 48$: 
$48 = 6 \cdot 8$. YES.

\medskip

For $N = 49$: 
$49 = 20 + 29$? No. 
$49 = 9 + 40 = 9 + 20 \cdot 2$. YES.

\medskip

For $N = 50$: 
$50 = 20 + 30 = 20 + 6 \cdot 5$. YES.

\textbf{Step 5: All $N \ge 50$ are achievable.} Once we reach 50 onwards, every number can be formed by adding packs of 6 (since we have 50, 51, 52, 53, 54, 55, and adding 6 cycles through all residues mod 6).

\textbf{Step 6: Maximum non-achievable.} The largest is $\boxed{43}$.

\paragraph{Example 17 (AMC 12 variant).} In how many ways can you pay exactly \$1.00 using only nickels (5 cents), dimes (10 cents), and quarters (25 cents)?

\paragraph{Solution.}
\textbf{Step 1: Set up the equation.} We need $5n + 10d + 25q = 100$ where $n, d, q \ge 0$ are integers.

Divide by 5: $n + 2d + 5q = 20$.

\textbf{Step 2: Fix $q$ and count $(n, d)$ pairs.} For each value of $q \in [0, 4]$:
\[
n + 2d = 20 - 5q
\]

For this to have non-negative solutions, we need $20 - 5q \ge 0$, so $q \le 4$.

\textbf{Step 3: Count solutions for each $q$.}
\begin{itemize}
    \item $q = 0$: $n + 2d = 20$. For each $d \in [0, 10]$, we get $n = 20 - 2d$. Count: 11 solutions.
    \item $q = 1$: $n + 2d = 15$. For each $d \in [0, 7]$, we get $n = 15 - 2d$. Count: 8 solutions.
    \item $q = 2$: $n + 2d = 10$. For each $d \in [0, 5]$, we get $n = 10 - 2d$. Count: 6 solutions.
    \item $q = 3$: $n + 2d = 5$. For each $d \in [0, 2]$, we get $n = 5 - 2d$. Count: 3 solutions.
    \item $q = 4$: $n + 2d = 0$. Only $d = 0, n = 0$. Count: 1 solution.
\end{itemize}

\textbf{Step 4: Total.}
\[
11 + 8 + 6 + 3 + 1 = 29
\]

Alternatively, using a generating function formula: the number of non-negative integer solutions to $n + 2d = m$ is $\lfloor m/2 \rfloor + 1$. So:
\begin{align*}
\sum_{q=0}^{4} (\lfloor (20-5q)/2 \rfloor + 1) &= 11 + 8 + 6 + 3 + 1 = 29
\end{align*}

\paragraph{Answer.} $\boxed{29}$

\section*{Closing Perspective}

Strong number theory is about \textbf{forcing outcomes}. Once primes, residues, and valuations are exposed, the answer is no longer a mystery — it becomes inevitable.

Master these techniques through repeated practice. The best number theorists recognize patterns instantly and choose the right tool without hesitation. Build this intuition by solving many problems and understanding why each technique works.

\end{document}
